\documentclass[UTF8, a4paper, 12pt]{ctexart}

% ==================== 1. 基础宏包 ====================
\usepackage{amsmath, amssymb, amsfonts} % 数学公式
\usepackage{graphicx}                 % 插图支持
\usepackage{booktabs}                 % 学术三线表
\usepackage{xcolor}                   % 颜色支持（代码高亮必需）
\usepackage{float}                    % 浮动体精确定位 [H]
\usepackage{geometry}                 % 页边距设置
\usepackage{listings}                 % 代码块排版
\usepackage{hyperref}                 % 超链接支持

% 页边距合理设置
\geometry{left=3cm, right=3cm, top=2.5cm, bottom=2.5cm}

\pagestyle{plain} % 页眉页脚样式

% ==================== 2. 工具配置 ====================
% 安全的listings代码块配置
\lstset{
    language=Python,
    basicstyle=\ttfamily\small,
    frame=single,
    keywordstyle=\color{blue},
    commentstyle=\color{gray},
    numbers=left,
    numberstyle=\tiny,
    breaklines=true,
    showstringspaces=false
}

% 修复版：安全的原生图片占位框（内部不打印含有潜在下划线的ID，彻底防止报错）
\newcommand{\figureplaceholder}[2]{
    \begin{figure}[H]
        \centering
        \framebox[10cm]{\rule{0pt}{4cm} \makebox[0pt]{[图片占位符: #2]}}
        \caption{#2}
        \label{fig:#1}
    \end{figure}
}

% ==================== 3. 超链接 (放在最后加载) ====================
\usepackage{hyperref}
\hypersetup{
    colorlinks=true,
    linkcolor=black,
    citecolor=black,
    urlcolor=blue
}

% ==================== 4. 论文元数据 ====================
\title{\textbf{此处输入您的论文题目}}
\author{作者姓名}
\date{\today}

% ==================== 5. 正文开始 ====================
\begin{document}

\maketitle

\begin{abstract}
此处输入您的中文摘要内容。摘要应当简要说明研究的背景、目的、方法、主要结果和结论。

\vspace{1em}
\noindent\textbf{关键词：} 关键词1；关键词2；关键词3
\end{abstract}

\newpage
\tableofcontents
\newpage

\section{问题重述}
\subsection{问题背景}
此处输入背景介绍...

\subsection{需要解决的问题}
\noindent\textbf{问题一}：具体描述...

\noindent\textbf{问题二}：具体描述...

\section{问题分析}
\subsection{关于问题一的分析}
分析内容...

\subsection{关于问题二的分析}
分析内容...

\subsection{关于问题三的分析}
分析内容...

\section{模型假设}
\begin{enumerate}
    \item 假设一...
    \item 假设二...
\end{enumerate}

\section{符号说明}
\begin{table}[H]
    \centering
    \caption{主要符号物理含义}
    \label{tab:symbols}
    \begin{tabular}{ccc}
        \toprule
        \textbf{符号} & \textbf{含义} & \textbf{单位} \\
        \midrule
        $d$ & 厚度 & m \\
        $n$ & 折射率 & / \\
        \bottomrule
    \end{tabular}
\end{table}

\section{模型的建立与求解}
\subsection{公式推导}
这里插入公式：
\begin{equation}\label{eq1}
    \begin{split}
    F = ma\\
    v = \frac{dx}{dt}
    \end{split}
\end{equation}
正如公式\autoref{eq1}所示，力与加速度的关系是线性的。
$$v=at$$
我们还有一个公式
\begin{equation}
    E = mc^2
\end{equation}

\subsection{算法与伪代码描述}
\begin{lstlisting}[language=Python]
# Write your algorithm or pseudocode here
def core_algorithm(input_data):
    if input_data is not None:
        return True
    return False
\end{lstlisting}

\subsection{结果分析}
\begin{figure}
    \centering
    \includegraphics[width=0.6\textwidth]{相对路径（斜杠方向换一下）} % 替换为实际图片路径
    \caption{示例图片说明}
    \label{fig1}
\end{figure}


\section{模型评价}
\subsection{优点}
第一点优点...
\subsection{缺点}
第一点缺点...

\begin{thebibliography}{99}
    \bibitem{ref1} 相关文献...
\end{thebibliography}

\newpage
\appendix
\section{附录代码}
\begin{lstlisting}[language=Python]
# Python code example
def calculate_thickness(wavelength, index):
    return wavelength / (2 * index)
\end{lstlisting}

\end{document}